\chapter{Prostate Biopsy}

\section*{Overview}
A prostate biopsy removes small tissue cores from the prostate gland using a spring-loaded needle, typically guided by transrectal ultrasound (TRUS) or MRI-fusion.

\section*{Alternative Names}
TRUS biopsy, Transrectal ultrasound biopsy

\section*{Principle}
An ultrasound probe in the rectum guides a spring-loaded biopsy needle into the prostate gland, obtaining multiple cores of tissue from different regions, particularly the peripheral zone, where most cancers arise. The cores are examined microscopically for cancer cells.
\section*{History}
Early biopsies were performed blindly. Transrectal ultrasound guidance was introduced in the 1980s, standardizing the 12-core systematic biopsy.

\section*{Why It Is Done}
Ordered when prostate cancer is suspected due to an elevated Prostate-Specific Antigen (PSA) level or an abnormal Digital Rectal Exam (DRE).

\section*{Is This a Primary or Secondary Test or Procedure}
Primary diagnostic test to confirm and grade prostate cancer.

\section*{How It Is Performed}
Under local anesthesia. An ultrasound probe is placed in the rectum to visualize the prostate. A biopsy needle is advanced through the rectal wall (transrectal) or perineum (transperineal) to take 12 tissue samples.

\section*{Preparation}
Antibiotics and bowel preparation depend on the biopsy route, local infection patterns, and clinician instructions. Tell the team about anticoagulants, antiplatelet medicines, allergies, and prior resistant infections. Do not stop blood-thinning medicines without an individualized plan from the prescribing and procedural teams.

\section*{Contraindications and Precautions}
Acute prostatitis, or uncorrected bleeding disorders.

\section*{Risks}
Hematuria, hematochezia (blood in stool), hematospermia (blood in semen, common for weeks), urinary retention, or severe sepsis (due to rectal bacteria translocation).

\section*{Aftercare and Recovery}
Complete the full course of antibiotics. Drink plenty of water. Avoid strenuous activity for 2-3 days.

\section*{Outcome and Follow-up}
Biopsy cores are evaluated under a microscope. Cancer is graded using the Gleason Score (e.g., Gleason 3+3=6 to 5+5=10).

\section*{Expected Results}
Benign prostatic tissue; no evidence of prostatic adenocarcinoma or high-grade PIN.

\section*{Follow-up}
Pathology review; discussion of active surveillance, surgery, or radiation if cancer is found.

\section*{When to Seek Medical Care}
Follow the procedure team's specific instructions. Seek urgent medical assessment for severe or worsening pain, uncontrolled bleeding, fever or signs of infection, trouble breathing, chest pain, fainting, new weakness or confusion, inability to urinate, or any rapidly worsening symptom. The appropriate warning signs depend on the procedure and the patient's medical condition.

\section*{References}
\begin{enumerate}
  \item MedlinePlus. Prostate Biopsy. U.S. National Library of Medicine; 2025.
  \item Current Medical Diagnosis and Treatment. McGraw-Hill; 2025.
  \item American Urological Association. Prostate needle biopsy patient-safety guidance. \url{https://www.auanet.org/documents/practices-resources/quality/patient_safety/prostate-needle-biopsy-white-paper.pdf} (accessed 2026-08-16).
\end{enumerate}

\clearpage
